\documentclass[11pt,a4paper,openany]{book}

\usepackage[T1]{fontenc}
\usepackage[utf8]{inputenc}
\usepackage{lmodern}
\usepackage{microtype}
\usepackage[a4paper,margin=1in]{geometry}
\usepackage{amsmath,amssymb,amsthm,mathtools,bm}
\usepackage{booktabs,array,longtable,tabularx}
\usepackage{enumitem}
\usepackage{xcolor}
\usepackage{hyperref}
\usepackage{fancyhdr}
\usepackage{titlesec}
\usepackage{listings}
\usepackage[most]{tcolorbox}
\usepackage{caption}

\definecolor{eldaBlue}{HTML}{16324F}
\definecolor{eldaGold}{HTML}{B5872C}
\definecolor{eldaSlate}{HTML}{374151}
\definecolor{eldaLight}{HTML}{F5F7FA}
\definecolor{eldaCode}{HTML}{F7F7F8}
\definecolor{eldaGreen}{HTML}{2F6B53}
\definecolor{eldaRed}{HTML}{8B2E2E}

\hypersetup{
    colorlinks=true,
    linkcolor=eldaBlue,
    urlcolor=eldaBlue,
    citecolor=eldaBlue,
    pdftitle={ELDA Linear Algebra Library Booklet},
    pdfauthor={Codex for the ELDA project}
}

\pdfstringdefDisableCommands{%
  \def\code#1{#1}%
  \def\ELDA{ELDA}%
}

\setcounter{tocdepth}{2}
\setcounter{secnumdepth}{3}
\setlength{\headheight}{14pt}

\newcommand{\ELDA}{\textsc{ELDA}}
\DeclareRobustCommand{\code}[1]{\texttt{\detokenize{#1}}}
\newcommand{\R}{\mathbb{R}}
\DeclareMathOperator{\rankop}{rank}
\DeclareMathOperator{\tr}{tr}
\DeclareMathOperator{\adj}{adj}

\titleformat{\chapter}[display]
  {\normalfont\bfseries\color{eldaBlue}}
  {\filleft\Huge\thechapter}
  {1ex}
  {\Huge}
  [\vspace{1ex}\titlerule]

\titleformat{\section}
  {\Large\bfseries\color{eldaBlue}}
  {\thesection}{0.75em}{}

\titleformat{\subsection}
  {\large\bfseries\color{eldaSlate}}
  {\thesubsection}{0.75em}{}

\pagestyle{fancy}
\fancyhf{}
\fancyhead[LE,RO]{\color{eldaSlate}\thepage}
\fancyhead[LO]{\color{eldaSlate}\nouppercase{\rightmark}}
\fancyhead[RE]{\color{eldaSlate}\nouppercase{\leftmark}}
\renewcommand{\headrulewidth}{0.4pt}
\renewcommand{\footrulewidth}{0pt}

\lstdefinestyle{eldaCpp}{
    language=C++,
    basicstyle=\ttfamily\small,
    keywordstyle=\color{eldaBlue}\bfseries,
    commentstyle=\color{eldaGreen},
    stringstyle=\color{eldaRed},
    showstringspaces=false,
    frame=single,
    rulecolor=\color{black!20},
    backgroundcolor=\color{eldaCode},
    columns=fullflexible,
    keepspaces=true
}

\newtcolorbox{mathbox}[1][]{
    enhanced,
    breakable,
    colback=eldaLight,
    colframe=eldaBlue,
    boxrule=0.8pt,
    left=8pt,
    right=8pt,
    top=8pt,
    bottom=8pt,
    title=#1,
    fonttitle=\bfseries
}

\newtcolorbox{examplebox}[1][]{
    enhanced,
    breakable,
    colback=white,
    colframe=eldaGold,
    boxrule=0.8pt,
    left=8pt,
    right=8pt,
    top=8pt,
    bottom=8pt,
    title=#1,
    fonttitle=\bfseries
}

\newtcolorbox{techbox}[1][]{
    enhanced,
    breakable,
    colback=eldaCode,
    colframe=eldaSlate,
    boxrule=0.8pt,
    left=8pt,
    right=8pt,
    top=8pt,
    bottom=8pt,
    title=#1,
    fonttitle=\bfseries
}

\begin{document}

\begin{titlepage}
    \centering
    \vspace*{2cm}
    {\Huge\bfseries\color{eldaBlue} ELDA \\[0.35cm]}
    {\LARGE\bfseries Linear Algebra Library Booklet \\[0.8cm]}
    {\large Mathematical reference, engineering guide, and practical starter manual\par}
    \vfill
    \rule{0.82\textwidth}{1.4pt}\\[0.5cm]
    {\Large\color{eldaSlate} Dense matrix operations, elimination algorithms, system solving, \\[0.15cm]
    norms, characteristic polynomials, and homogeneous transforms\par}
    \rule{0.82\textwidth}{1.4pt}
    \vfill
    \begin{minipage}{0.82\textwidth}
        \centering
        \begin{techbox}[Document Scope]
            This booklet documents the current ELDA C++17 library found under \code{include/elda/} and \code{src/}. It explains the mathematics used by the library, the engineering structure of the codebase, and the intended usage of every major public API entry point.
        \end{techbox}
    \end{minipage}
    \vfill
    {\large Prepared for the project repository}\\[0.2cm]
    {\large \today}
\end{titlepage}

\frontmatter

\chapter*{Preface}
\addcontentsline{toc}{chapter}{Preface}
\ELDA{} is a compact linear algebra library written in modern C++ and organized as a CMake project. The library exposes a dense matrix type, arithmetic helpers, elimination-based routines, structural invariants such as determinant and rank, and 2D/3D transformation generators based on homogeneous coordinates.

This booklet serves three purposes:
\begin{enumerate}[label=\arabic*.]
    \item It acts as a starter guide for someone who wants to build and use the library quickly.
    \item It explains the mathematics that justify each public function, so the code is tied back to linear algebra rather than treated as a black box.
    \item It records technical structure and implementation behavior so future maintenance work can stay coherent.
\end{enumerate}

\tableofcontents

\mainmatter

\chapter{Library Overview and Starter Guide}

\section{What ELDA provides}
At its core, \ELDA{} provides a \code{linalg::matrix} type whose entries are stored as \code{double} values in a row-major \code{std::vector<std::vector<double>>}. The public API is grouped into three headers:
\begin{itemize}
    \item \code{elda/matrix.hpp}: the matrix type and core algebraic routines,
    \item \code{elda/transforms.hpp}: 2D and 3D homogeneous transformation matrices,
    \item \code{elda/vector_utils.hpp}: one library-specific comparison helper.
\end{itemize}

\begin{techbox}[Project Layout]
\begin{tabularx}{\linewidth}{@{}lX@{}}
\toprule
Path & Purpose \\
\midrule
\code{CMakeLists.txt} & Defines the static library target \code{elda} and the demo executable \code{main}. \\
\code{include/elda/matrix.hpp} & Public declaration of the matrix type and most free functions. \\
\code{include/elda/transforms.hpp} & Public declaration of 2D and 3D transform helpers. \\
\code{include/elda/vector_utils.hpp} & Public declaration of \code{check_similar}. \\
\code{include/elda/linalg.hpp} & Umbrella header that includes the whole public API. \\
\code{src/matrix.cpp} & Implementation of matrix arithmetic, elimination, determinant, inverse, and spectral helpers. \\
\code{src/transforms.cpp} & Implementation of translation, scaling, and rotation matrices. \\
\code{src/vector_utils.cpp} & Implementation of the zero-row filtered comparison utility. \\
\code{main.cpp} & Small demo program that constructs a matrix, runs Gaussian elimination, and prints the result. \\
\bottomrule
\end{tabularx}
\end{techbox}

\section{Build the project}
The library is configured with CMake and targets C++17.

\begin{lstlisting}[style=eldaCpp,caption={Build commands}]
cmake -S . -B build
cmake --build build
\end{lstlisting}

After a successful build, the main artifacts are:
\begin{itemize}
    \item \code{build/libelda.a}: the static library,
    \item \code{build/main}: the demo executable.
\end{itemize}

\section{Minimal usage example}
The fastest way to start is to include the umbrella header and work with a small dense matrix.

\begin{lstlisting}[style=eldaCpp,caption={Minimal ELDA program}]
#include <elda/linalg.hpp>
#include <iostream>

using namespace linalg;

int main() {
    matrix a(2, 2);
    a.arr = {{1, 2}, {3, 4}};

    matrix b = identity(2);
    matrix c = a + b;

    c.print();
}
\end{lstlisting}

\begin{examplebox}[Starter interpretation]
The matrix
\[
A = \begin{bmatrix}
1 & 2 \\
3 & 4
\end{bmatrix}
\]
is added to the identity matrix
\[
I_2 = \begin{bmatrix}
1 & 0 \\
0 & 1
\end{bmatrix},
\]
which yields
\[
A + I_2 =
\begin{bmatrix}
2 & 2 \\
3 & 5
\end{bmatrix}.
\]
This is entrywise matrix addition, so each diagonal entry increases by one.
\end{examplebox}

\section{Who this booklet is for}
This document is written for three audiences:
\begin{itemize}
    \item users who want to apply the library correctly,
    \item maintainers who need to understand why an algorithm exists,
    \item learners who want the code tied back to linear algebra concepts.
\end{itemize}

\chapter{Mathematical and Technical Foundations}

\section{Notation used in this booklet}
Unless stated otherwise, a matrix
\[
A = (a_{ij}) \in \R^{m \times n}
\]
has \(m\) rows and \(n\) columns. In the C++ code, the dimensions are stored in the public members \code{row} and \code{col}, and indexing is zero-based because the underlying storage is a C++ vector.

\begin{mathbox}[Indexing convention]
Mathematics often writes \(a_{11}\) for the first element of a matrix. ELDA stores that same element at \code{arr[0][0]}. In formulas below we use the standard mathematical notation, but examples mention the corresponding C++ indices whenever confusion is possible.
\end{mathbox}

\section{Dense storage model}
The matrix class stores data as
\[
\texttt{std::vector<std::vector<double>>}.
\]
This is easy to inspect, easy to print, and convenient for educational code. The tradeoff is that it is not as cache-friendly as a flat contiguous buffer. For the scale of a teaching-oriented or small library, the design is reasonable and transparent.

\section{Mutating versus non-mutating operations}
\ELDA{} has two styles of API:
\begin{itemize}
    \item \textbf{return-a-copy helpers}: arithmetic operators, row and column operations, transpose, inverse, adjoint, and transform factories return new matrices;
    \item \textbf{in-place reducers}: \code{echelon()}, \code{gaussian()}, \code{gauss_jordan()}, and \code{canonical()} mutate the matrix object and return an integer that records the sign change induced by row swaps.
\end{itemize}

\section{Core linear algebra ideas}
\subsection{Elementary row operations}
Three row operations preserve the solution set of a linear system:
\begin{enumerate}[label=\arabic*.]
    \item swap two rows,
    \item multiply a row by a nonzero scalar,
    \item replace one row by itself plus a multiple of another row.
\end{enumerate}

These operations are the engine behind Gaussian elimination, rank computation, determinant reduction, and matrix inversion. ELDA also exposes column analogues of the same operations.

\subsection{Homogeneous coordinates}
The transformation helpers in \code{transforms.hpp} use homogeneous coordinates. A 2D point \((x, y)\) is encoded as
\[
\begin{bmatrix}
x \\ y \\ 1
\end{bmatrix},
\]
and a 3D point \((x, y, z)\) is encoded as
\[
\begin{bmatrix}
x \\ y \\ z \\ 1
\end{bmatrix}.
\]
This lets translation, scaling, and rotation all be represented as matrix multiplication.

\chapter{Matrix Type, Constructors, and Basic API}

\section{Constructors}

\subsection{\code{matrix()}}
The default constructor builds a \(3 \times 3\) zero matrix.

\begin{mathbox}[Mathematical meaning]
The zero matrix \(0_{3 \times 3}\) is the additive identity in the vector space of \(3 \times 3\) real matrices. For every \(A \in \R^{3 \times 3}\),
\[
A + 0_{3 \times 3} = A.
\]
\end{mathbox}

\begin{lstlisting}[style=eldaCpp,caption={Default constructor example}]
matrix z;
z.print();
\end{lstlisting}

\begin{examplebox}[Example]
The code above prints a \(3 \times 3\) grid of zeros.
\end{examplebox}

\subsection{\code{matrix(int r, int c)}}
This constructor builds the zero matrix \(0_{r \times c}\).

\begin{lstlisting}[style=eldaCpp,caption={Sized constructor example}]
matrix a(2, 3);
a.arr = {{1, 2, 3}, {4, 5, 6}};
\end{lstlisting}

\begin{examplebox}[Example]
Before assigning entries, \code{a} is
\[
\begin{bmatrix}
0 & 0 & 0 \\
0 & 0 & 0
\end{bmatrix}.
\]
\end{examplebox}

\section{Element access and I/O}

\subsection{\code{get_element(i, j)}}
This returns the entry \(a_{ij}\) with no bounds check. The mathematics is direct: it selects one coordinate of the matrix.

\begin{examplebox}[Example]
If
\[
A = \begin{bmatrix}
7 & 8 \\
9 & 10
\end{bmatrix},
\]
then \code{get_element(1, 0)} returns \(9\), because C++ index \((1,0)\) is the second row, first column.
\end{examplebox}

\subsection{\code{ref_element(i, j)}}
This returns a mutable pointer to the selected element. In mathematical terms it gives write access to one coordinate of the matrix.

\begin{lstlisting}[style=eldaCpp,caption={Pointer-based update}]
matrix a(2, 2);
*a.ref_element(0, 1) = 5.0;
\end{lstlisting}

\subsection{\code{input()}}
This reads all entries row by row from standard input. Mathematically it populates a matrix by specifying the coordinates \(a_{ij}\) in row-major order.

\subsection{\code{print()}}
This writes the matrix one row per line. It is a presentation helper rather than a mathematical operation.

\section{Assignment and comparisons}

\subsection{\code{operator=(matrix m2)}}
The library defines assignment as a shape-preserving copy. If the dimensions of the left-hand side and right-hand side differ, an exception is thrown.

\begin{techbox}[Important implementation note]
Standard C++ assignment usually resizes the destination object when needed. ELDA's current assignment operator instead requires equal dimensions. That is a library design decision, not a language requirement.
\end{techbox}

\begin{lstlisting}[style=eldaCpp,caption={Assignment example}]
matrix a(2, 2), b(2, 2);
b.arr = {{1, 0}, {0, 1}};
a = b;
\end{lstlisting}

\begin{examplebox}[Example]
After assignment, \(A = B = I_2\).
\end{examplebox}

\subsection{\code{operator==(matrix m1, matrix m2)}}
This function checks exact entrywise equality:
\[
A = B \quad \Longleftrightarrow \quad a_{ij} = b_{ij} \text{ for all } i,j.
\]

\begin{examplebox}[Example]
\[
\begin{bmatrix}
1 & 2 \\
3 & 4
\end{bmatrix}
=
\begin{bmatrix}
1 & 2 \\
3 & 4
\end{bmatrix},
\]
so the comparison returns \code{true}.
\end{examplebox}

\subsection{\code{shape_comp(matrix m1, matrix m2)}}
This tests only the dimensions. It corresponds to the statement \(A, B \in \R^{m \times n}\) for the same \(m,n\).

\begin{examplebox}[Example]
A \(2 \times 3\) matrix and another \(2 \times 3\) matrix are shape-compatible even if their entries are different.
\end{examplebox}

\section{Identity and structure helpers}

\subsection{\code{identity(int n)}}
This builds the identity matrix
\[
I_n =
\begin{bmatrix}
1 & 0 & \cdots & 0 \\
0 & 1 & \cdots & 0 \\
\vdots & \vdots & \ddots & \vdots \\
0 & 0 & \cdots & 1
\end{bmatrix}.
\]

\begin{mathbox}[Why it matters]
The identity matrix is the multiplicative identity:
\[
AI_n = A, \qquad I_nA = A.
\]
It is central to inversion, matrix powers, and transform construction.
\end{mathbox}

\begin{examplebox}[Example]
\code{identity(3)} returns \(I_3\), the matrix with ones on the main diagonal and zeros elsewhere.
\end{examplebox}

\subsection{\code{get_row(int r)}}
This returns a matrix with row \(r\) preserved and all other rows set to zero. If \(A\) is \(m \times n\), the result is also \(m \times n\), not a \(1 \times n\) row vector.

\begin{examplebox}[Example]
If
\[
A = \begin{bmatrix}
1 & 2 \\
3 & 4
\end{bmatrix},
\]
then \code{get_row(1)} returns
\[
\begin{bmatrix}
0 & 0 \\
3 & 4
\end{bmatrix}.
\]
\end{examplebox}

\subsection{\code{get_col(int c)}}
This returns a matrix with column \(c\) preserved and all other columns set to zero.

\begin{examplebox}[Example]
For the same matrix \(A\),
\[
\code{get_col(0)} =
\begin{bmatrix}
1 & 0 \\
3 & 0
\end{bmatrix}.
\]
\end{examplebox}

\chapter{Arithmetic Operations}

\section{\code{operator+(matrix m2)}}
Matrix addition is defined entrywise:
\[
(A+B)_{ij} = a_{ij} + b_{ij}.
\]
The operation is valid only when both matrices have the same dimensions.

\begin{examplebox}[Example]
\[
\begin{bmatrix}
1 & 2 \\
3 & 4
\end{bmatrix}
+
\begin{bmatrix}
5 & 6 \\
7 & 8
\end{bmatrix}
=
\begin{bmatrix}
6 & 8 \\
10 & 12
\end{bmatrix}.
\]
\end{examplebox}

\section{\code{operator-(matrix m2)}}
Subtraction is also entrywise:
\[
(A-B)_{ij} = a_{ij} - b_{ij}.
\]

\begin{examplebox}[Example]
\[
\begin{bmatrix}
5 & 6 \\
7 & 8
\end{bmatrix}
-
\begin{bmatrix}
1 & 2 \\
3 & 4
\end{bmatrix}
=
\begin{bmatrix}
4 & 4 \\
4 & 4
\end{bmatrix}.
\]
\end{examplebox}

\section{\code{operator*(double d)}}
Scalar multiplication rescales every entry:
\[
(dA)_{ij} = d \cdot a_{ij}.
\]

\begin{examplebox}[Example]
\[
2
\begin{bmatrix}
1 & -1 \\
3 & 0
\end{bmatrix}
=
\begin{bmatrix}
2 & -2 \\
6 & 0
\end{bmatrix}.
\]
In ELDA this is written as \code{A * 2.0}.
\end{examplebox}

\section{\code{operator*(matrix m2)}}
Matrix multiplication is defined by row-column pairing:
\[
(AB)_{ij} = \sum_{k=1}^{n} a_{ik} b_{kj},
\]
where \(A \in \R^{m \times n}\) and \(B \in \R^{n \times p}\). The result lies in \(\R^{m \times p}\).

\begin{mathbox}[Interpretation]
The entry \((i,j)\) of \(AB\) is the dot product of row \(i\) of \(A\) with column \(j\) of \(B\). This is why the inner dimensions must match.
\end{mathbox}

\begin{examplebox}[Example]
\[
A =
\begin{bmatrix}
1 & 2 \\
3 & 4
\end{bmatrix},
\qquad
B =
\begin{bmatrix}
2 & 0 \\
1 & 5
\end{bmatrix}.
\]
Then
\[
AB =
\begin{bmatrix}
1\cdot2 + 2\cdot1 & 1\cdot0 + 2\cdot5 \\
3\cdot2 + 4\cdot1 & 3\cdot0 + 4\cdot5
\end{bmatrix}
=
\begin{bmatrix}
4 & 10 \\
10 & 20
\end{bmatrix}.
\]
\end{examplebox}

\chapter{Elementary Row and Column Operations}

\section{\code{row_op(int i, double coeff, int j)}}
This performs the elementary row replacement
\[
R_i \leftarrow R_i + \text{coeff} \cdot R_j.
\]

\begin{mathbox}[Why this matters]
Row replacement preserves the solution set of a linear system and preserves the determinant. It is the workhorse operation in elimination.
\end{mathbox}

\begin{examplebox}[Example]
Start with
\[
\begin{bmatrix}
1 & 2 \\
3 & 4
\end{bmatrix}.
\]
Applying \code{row_op(1, -3, 0)} means
\[
R_2 \leftarrow R_2 - 3R_1,
\]
so the result is
\[
\begin{bmatrix}
1 & 2 \\
0 & -2
\end{bmatrix}.
\]
\end{examplebox}

\section{\code{col_op(int i, double coeff, int j)}}
This is the column analogue:
\[
C_i \leftarrow C_i + \text{coeff} \cdot C_j.
\]
Column operations are useful in canonical forms and matrix equivalence arguments.

\begin{examplebox}[Example]
For
\[
\begin{bmatrix}
1 & 2 \\
3 & 4
\end{bmatrix},
\]
applying \code{col_op(1, -2, 0)} gives
\[
C_2 \leftarrow C_2 - 2C_1,
\qquad
\begin{bmatrix}
1 & 0 \\
3 & -2
\end{bmatrix}.
\]
\end{examplebox}

\section{\code{row_swap(int i, int j)} and \code{col_swap(int i, int j)}}
Swapping rows or columns permutes the matrix.

\begin{mathbox}[Determinant effect]
Swapping two rows multiplies the determinant by \(-1\). Swapping two columns does the same. This sign flip is why the elimination routines keep track of swap parity.
\end{mathbox}

\begin{examplebox}[Example]
\[
\begin{bmatrix}
1 & 2 \\
3 & 4
\end{bmatrix}
\xrightarrow{\text{swap rows}}
\begin{bmatrix}
3 & 4 \\
1 & 2
\end{bmatrix}.
\]
\end{examplebox}

\section{\code{row_multi(int i, double factor)} and \code{col_multi(int i, double factor)}}
These scale one row or one column:
\[
R_i \leftarrow \lambda R_i,
\qquad
C_j \leftarrow \lambda C_j.
\]

\begin{mathbox}[Determinant effect]
If a row of a square matrix is multiplied by \(\lambda\), then the determinant is multiplied by \(\lambda\). Gaussian elimination uses row scaling to normalize pivots.
\end{mathbox}

\begin{examplebox}[Example]
\[
\begin{bmatrix}
1 & 2 \\
3 & 4
\end{bmatrix}
\xrightarrow{R_1 \leftarrow 2R_1}
\begin{bmatrix}
2 & 4 \\
3 & 4
\end{bmatrix}.
\]
\end{examplebox}

\chapter{Elimination, Rank, and Linear Systems}

\section{\code{echelon()}}
This mutates the matrix into row-echelon form by searching for a nonzero pivot in each column, swapping that row upward if necessary, and eliminating entries below the pivot.

\begin{mathbox}[Row-echelon form]
A matrix is in row-echelon form when:
\begin{enumerate}[label=\arabic*.]
    \item every nonzero row starts with a leading nonzero entry called a pivot,
    \item pivot positions move strictly to the right as you move downward,
    \item entries below each pivot are zero.
\end{enumerate}
\end{mathbox}

\begin{examplebox}[Example]
Let
\[
A =
\begin{bmatrix}
2 & 1 \\
4 & 3
\end{bmatrix}.
\]
After \code{echelon()}, ELDA performs
\[
R_2 \leftarrow R_2 - 2R_1,
\]
which produces
\[
\begin{bmatrix}
2 & 1 \\
0 & 1
\end{bmatrix}.
\]
This is already in row-echelon form.
\end{examplebox}

\section{\code{gaussian()}}
This is a normalized Gaussian elimination. It first creates a pivot, then scales the pivot row so the pivot becomes \(1\), and finally clears entries below it.

\begin{mathbox}[Gaussian elimination with normalized pivots]
For pivot position \((k,k)\), ELDA applies
\[
R_k \leftarrow \frac{1}{a_{kk}}R_k
\]
when \(a_{kk} \neq 0\), then uses row replacement to zero all entries below the pivot. The resulting form is upper row-echelon with leading ones whenever possible.
\end{mathbox}

\begin{examplebox}[Example]
For
\[
\begin{bmatrix}
3 & 6 \\
6 & 15
\end{bmatrix},
\]
the first pivot row is scaled by \(1/3\), giving
\[
\begin{bmatrix}
1 & 2 \\
6 & 15
\end{bmatrix},
\]
then \(R_2 \leftarrow R_2 - 6R_1\), yielding
\[
\begin{bmatrix}
1 & 2 \\
0 & 3
\end{bmatrix},
\]
and the second pivot row is scaled to obtain
\[
\begin{bmatrix}
1 & 2 \\
0 & 1
\end{bmatrix}.
\]
\end{examplebox}

\section{\code{gauss_jordan()}}
This calls \code{gaussian()} first, then eliminates the entries above each pivot to obtain reduced row-echelon form.

\begin{mathbox}[Reduced row-echelon form]
A matrix is in reduced row-echelon form when it is in row-echelon form and every pivot column contains a single \(1\) with zeros everywhere else.
\end{mathbox}

\begin{examplebox}[Example]
Starting from
\[
\begin{bmatrix}
1 & 2 \\
0 & 1
\end{bmatrix},
\]
the routine clears the entry above the second pivot:
\[
R_1 \leftarrow R_1 - 2R_2,
\]
producing
\[
\begin{bmatrix}
1 & 0 \\
0 & 1
\end{bmatrix}.
\]
\end{examplebox}

\section{\code{canonical()}}
This routine first runs normalized Gaussian elimination, then applies column operations to clear entries to the right of each pivot. The effect is a mixed row-column reduction that pushes the matrix toward a diagonal or canonical pattern.

\begin{mathbox}[Conceptual role]
Row operations preserve the solution set of a linear system, while column operations change the coordinate representation of the domain. Combining them is natural when studying matrix equivalence or canonical forms, rather than only solving systems.
\end{mathbox}

\begin{examplebox}[Example]
If Gaussian elimination has already produced
\[
\begin{bmatrix}
1 & 2 \\
0 & 1
\end{bmatrix},
\]
then a column operation
\[
C_2 \leftarrow C_2 - 2C_1
\]
transforms it into the identity matrix. This is the kind of diagonal-cleaning behavior that \code{canonical()} performs.
\end{examplebox}

\section{\code{rank()}}
The rank of a matrix is the dimension of its row space, equivalently the number of pivots in row-echelon form:
\[
\rankop(A) = \text{number of nonzero rows in an echelon form of } A.
\]
ELDA computes this by applying \code{gaussian()} to a copy and counting nonzero rows.

\begin{examplebox}[Example]
For
\[
\begin{bmatrix}
1 & 2 \\
2 & 4
\end{bmatrix},
\]
the second row is a multiple of the first, so there is only one pivot. Hence \(\rankop(A)=1\).
\end{examplebox}

\section{\code{free_variable()}}
For a square coefficient matrix \(A \in \R^{n \times n}\), this returns
\[
n - \rankop(A).
\]
This matches the nullity of \(A\) for square systems, by the rank-nullity theorem.

\begin{examplebox}[Example]
If a \(3 \times 3\) matrix has rank \(2\), then the routine returns \(3-2=1\), meaning one free variable in the homogeneous system \(Ax=0\).
\end{examplebox}

\section{\code{solve()}}
This function expects an augmented matrix of size \(n \times (n+1)\):
\[
\left[\begin{array}{c|c}
A & b
\end{array}\right].
\]
It runs normalized Gaussian elimination and then performs back substitution to produce the solution vector.

\begin{mathbox}[Back substitution]
If the reduced system is upper triangular,
\[
\begin{aligned}
x_n &= b_n, \\
x_{n-1} &= b_{n-1} - a_{n-1,n}x_n, \\
x_{n-2} &= b_{n-2} - a_{n-2,n-1}x_{n-1} - a_{n-2,n}x_n,
\end{aligned}
\]
and so on upward.
\end{mathbox}

\begin{examplebox}[Example]
Solve
\[
\begin{cases}
x + y = 3, \\
2x + 3y = 8.
\end{cases}
\]
The augmented matrix is
\[
\left[
\begin{array}{cc|c}
1 & 1 & 3 \\
2 & 3 & 8
\end{array}
\right].
\]
Gaussian elimination gives
\[
\left[
\begin{array}{cc|c}
1 & 1 & 3 \\
0 & 1 & 2
\end{array}
\right].
\]
Back substitution yields \(y=2\) and \(x=1\).
\end{examplebox}

\chapter{Determinants, Cofactors, and Inverses}

\section{\code{trace()}}
The trace is the sum of diagonal entries:
\[
\tr(A) = \sum_{i=1}^{\min(m,n)} a_{ii}.
\]
For square matrices, the trace equals the sum of eigenvalues counted with algebraic multiplicity.

\begin{examplebox}[Example]
\[
\tr\left(
\begin{bmatrix}
2 & 1 \\
5 & 3
\end{bmatrix}
\right) = 2 + 3 = 5.
\]
\end{examplebox}

\section{\code{diag_product()}}
This returns the product of the diagonal entries:
\[
\prod_{i=1}^{\min(m,n)} a_{ii}.
\]
For triangular square matrices, this product is exactly the determinant.

\begin{examplebox}[Example]
For
\[
\begin{bmatrix}
2 & 7 \\
0 & 5
\end{bmatrix},
\]
\code{diag_product()} returns \(2 \cdot 5 = 10\).
\end{examplebox}

\section{\code{det()}}
The determinant of a square matrix is computed by converting a copy of the matrix to row-echelon form and then multiplying the diagonal entries, adjusted by the sign of row swaps.

\begin{mathbox}[Why elimination works]
Two determinant facts make the algorithm valid:
\begin{enumerate}[label=\arabic*.]
    \item replacing one row by itself plus a multiple of another row leaves the determinant unchanged,
    \item swapping two rows changes the sign of the determinant.
\end{enumerate}
Once the matrix is upper triangular,
\[
\det(A) = \prod_i u_{ii}.
\]
\end{mathbox}

\begin{examplebox}[Example]
For
\[
A =
\begin{bmatrix}
1 & 2 \\
3 & 4
\end{bmatrix},
\]
perform
\[
R_2 \leftarrow R_2 - 3R_1
\]
to get
\[
\begin{bmatrix}
1 & 2 \\
0 & -2
\end{bmatrix}.
\]
No row swap occurred, so
\[
\det(A) = 1 \cdot (-2) = -2.
\]
\end{examplebox}

\section{\code{minor(int x, int y)}}
The minor of entry \((x,y)\) is the determinant of the submatrix obtained by deleting row \(x\) and column \(y\).

\begin{examplebox}[Example]
For
\[
\begin{bmatrix}
1 & 2 & 3 \\
4 & 5 & 6 \\
7 & 8 & 9
\end{bmatrix},
\]
the minor of the center entry is
\[
\det
\begin{bmatrix}
1 & 3 \\
7 & 9
\end{bmatrix}
= 9 - 21 = -12.
\]
\end{examplebox}

\section{\code{cofactor(int x, int y)}}
The cofactor adds the alternating sign pattern:
\[
C_{xy} = (-1)^{x+y} M_{xy},
\]
where \(M_{xy}\) is the minor.

\begin{examplebox}[Example]
If the minor at \((1,1)\) equals \(-12\), then the cofactor is
\[
(-1)^{1+1}(-12) = -12.
\]
If the same minor were at \((1,2)\), the cofactor would flip sign.
\end{examplebox}

\section{\code{adjoint()}}
The classical adjoint, also called the adjugate, is the transpose of the cofactor matrix:
\[
\adj(A) = C^\top.
\]

\begin{mathbox}[Inverse connection]
For a nonsingular square matrix,
\[
A^{-1} = \frac{1}{\det(A)} \adj(A).
\]
ELDA does not use this identity to compute the inverse, but the formula explains why the adjoint is mathematically important.
\end{mathbox}

\begin{examplebox}[Example]
For
\[
A =
\begin{bmatrix}
1 & 2 \\
3 & 4
\end{bmatrix},
\]
the adjugate is
\[
\adj(A) =
\begin{bmatrix}
4 & -2 \\
-3 & 1
\end{bmatrix}.
\]
\end{examplebox}

\section{\code{transpose()}}
The transpose exchanges rows and columns:
\[
(A^\top)_{ij} = a_{ji}.
\]
It is fundamental for symmetric matrices, orthogonality checks, and the Frobenius inner product.

\begin{examplebox}[Example]
\[
\begin{bmatrix}
1 & 2 & 3 \\
4 & 5 & 6
\end{bmatrix}^\top
=
\begin{bmatrix}
1 & 4 \\
2 & 5 \\
3 & 6
\end{bmatrix}.
\]
\end{examplebox}

\section{\code{inverse()}}
The inverse \(A^{-1}\) of a square matrix satisfies
\[
AA^{-1} = A^{-1}A = I.
\]
ELDA computes it through Gauss-Jordan elimination on the augmented block matrix
\[
\left[\begin{array}{c|c}
A & I
\end{array}\right].
\]
When the left block is reduced to \(I\), the right block becomes \(A^{-1}\).

\begin{examplebox}[Example]
For
\[
A =
\begin{bmatrix}
1 & 2 \\
3 & 4
\end{bmatrix},
\]
augment with the identity:
\[
\left[
\begin{array}{cc|cc}
1 & 2 & 1 & 0 \\
3 & 4 & 0 & 1
\end{array}
\right].
\]
After Gauss-Jordan elimination, the result is
\[
\left[
\begin{array}{cc|cc}
1 & 0 & -2 & 1 \\
0 & 1 & \frac{3}{2} & -\frac{1}{2}
\end{array}
\right],
\]
so
\[
A^{-1} =
\begin{bmatrix}
-2 & 1 \\
\frac{3}{2} & -\frac{1}{2}
\end{bmatrix}.
\]
\end{examplebox}

\chapter{Norms, Inner Products, Matrix Powers, and Characteristic Polynomials}

\section{\code{norm()}}
ELDA uses the Frobenius norm:
\[
\|A\|_F = \sqrt{\sum_{i,j} a_{ij}^2}.
\]
This is the Euclidean norm of the matrix entries viewed as one long vector.

\begin{examplebox}[Example]
For
\[
A =
\begin{bmatrix}
1 & 2 \\
2 & 1
\end{bmatrix},
\]
\[
\|A\|_F = \sqrt{1^2 + 2^2 + 2^2 + 1^2} = \sqrt{10}.
\]
\end{examplebox}

\section{\code{inner_product(matrix a, matrix b)}}
The function computes the Frobenius inner product:
\[
\langle A, B \rangle_F = \tr(A^\top B) = \sum_{i,j} a_{ij}b_{ij}.
\]

\begin{examplebox}[Example]
If
\[
A =
\begin{bmatrix}
1 & 2 \\
0 & 1
\end{bmatrix},
\qquad
B =
\begin{bmatrix}
3 & 4 \\
5 & 6
\end{bmatrix},
\]
then
\[
\langle A, B \rangle_F = 1\cdot3 + 2\cdot4 + 0\cdot5 + 1\cdot6 = 17.
\]
\end{examplebox}

\section{\code{angle(matrix a, matrix b)}}
The angle between two matrices is defined through the Frobenius inner product:
\[
\theta = \cos^{-1}\left(\frac{\langle A, B \rangle_F}{\|A\|_F \|B\|_F}\right).
\]
This generalizes the Euclidean angle between vectors.

\begin{examplebox}[Example]
If \(A = I_2\) and \(B = I_2\), then the numerator equals the product of the norms, so the cosine is \(1\) and the angle is \(0\).
\end{examplebox}

\section{\code{matpow(matrix mat, long long expo)}}
This computes \(A^k\) for a nonnegative integer \(k\) using repeated squaring.

\begin{mathbox}[Repeated squaring]
Instead of multiplying \(A\) by itself \(k\) times in sequence, repeated squaring uses the binary expansion of \(k\). For example,
\[
A^{13} = A^{8}A^{4}A.
\]
This reduces the number of matrix multiplications from \(O(k)\) to \(O(\log k)\).
\end{mathbox}

\begin{examplebox}[Example]
For
\[
A =
\begin{bmatrix}
1 & 1 \\
0 & 1
\end{bmatrix},
\]
one gets
\[
A^2 =
\begin{bmatrix}
1 & 2 \\
0 & 1
\end{bmatrix},
\qquad
A^3 =
\begin{bmatrix}
1 & 3 \\
0 & 1
\end{bmatrix}.
\]
\end{examplebox}

\section{\code{char_poly()}}
This routine implements the Faddeev-LeVerrier recurrence to compute the coefficients of the characteristic polynomial
\[
p_A(\lambda) = \det(\lambda I - A)
= \lambda^n + c_1\lambda^{n-1} + \cdots + c_n
\]
for an \(n \times n\) matrix.

\begin{mathbox}[Faddeev-LeVerrier recurrence]
The algorithm used by ELDA follows the pattern
\[
B_0 = I,
\qquad
c_k = -\frac{1}{k}\tr(AB_{k-1}),
\qquad
B_k = AB_{k-1} + c_k I.
\]
Each \(c_k\) becomes one coefficient of the characteristic polynomial.
\end{mathbox}

\begin{examplebox}[Example]
For
\[
A =
\begin{bmatrix}
1 & 2 \\
3 & 4
\end{bmatrix},
\]
the characteristic polynomial is
\[
\lambda^2 - 5\lambda - 2.
\]
So the returned row vector is
\[
\begin{bmatrix}
1 & -5 & -2
\end{bmatrix}.
\]
\end{examplebox}

\section{\code{check_ortho(matrix mat)}}
An orthogonal real matrix satisfies
\[
A^\top A = I,
\]
which is equivalent to
\[
A^{-1} = A^\top.
\]
ELDA checks the second formulation by comparing the transpose and inverse directly.

\begin{examplebox}[Example]
The planar rotation matrix
\[
R(\theta) =
\begin{bmatrix}
\cos\theta & -\sin\theta \\
\sin\theta & \cos\theta
\end{bmatrix}
\]
is orthogonal for all real \(\theta\).
\end{examplebox}

\section{\code{check_unitary(matrix mat)}}
In standard linear algebra, a complex matrix is unitary when
\[
U^\ast U = I,
\]
where \(U^\ast\) is the conjugate transpose. ELDA's current helper compares \code{transpose()} and \code{adjoint()}, where \code{adjoint()} means the classical adjugate, not the conjugate transpose.

\begin{techbox}[Important terminology note]
This means \code{check_unitary()} should be understood as a legacy project helper rather than a standard mathematical unitary test.
\end{techbox}

\begin{examplebox}[Example]
For the identity matrix, the transpose and adjugate are both equal to the identity, so the helper returns \code{true}. That does not validate the helper as a complete unitary predicate in the usual mathematical sense.
\end{examplebox}

\chapter{Transformation Helpers}

\section{2D transformation matrices}
\subsection{\code{shift(double dx, double dy)}}
This returns the homogeneous translation matrix
\[
T(dx,dy) =
\begin{bmatrix}
1 & 0 & dx \\
0 & 1 & dy \\
0 & 0 & 1
\end{bmatrix}.
\]

\begin{examplebox}[Example]
Applying this matrix to the homogeneous point \((x,y,1)^\top\) gives
\[
\begin{bmatrix}
1 & 0 & dx \\
0 & 1 & dy \\
0 & 0 & 1
\end{bmatrix}
\begin{bmatrix}
x \\ y \\ 1
\end{bmatrix}
=
\begin{bmatrix}
x + dx \\ y + dy \\ 1
\end{bmatrix}.
\]
\end{examplebox}

\subsection{\code{scale(double kx, double ky)}}
This returns
\[
S(k_x,k_y) =
\begin{bmatrix}
k_x & 0 & 0 \\
0 & k_y & 0 \\
0 & 0 & 1
\end{bmatrix},
\]
which scales the \(x\)-coordinate by \(k_x\) and the \(y\)-coordinate by \(k_y\).

\begin{examplebox}[Example]
If \((x,y)=(2,3)\) and \((k_x,k_y)=(4,5)\), then the scaled point is \((8,15)\).
\end{examplebox}

\subsection{\code{rotate(double angle)}}
This returns the standard planar rotation matrix in homogeneous form:
\[
R(\theta) =
\begin{bmatrix}
\cos\theta & -\sin\theta & 0 \\
\sin\theta & \cos\theta & 0 \\
0 & 0 & 1
\end{bmatrix}.
\]
Angles are measured in radians. The header also defines the convenience constant \code{PI = 3.141593}.

\begin{examplebox}[Example]
When \(\theta=\pi/2\),
\[
R(\pi/2)
\begin{bmatrix}
1 \\ 0 \\ 1
\end{bmatrix}
=
\begin{bmatrix}
0 \\ 1 \\ 1
\end{bmatrix},
\]
so the point \((1,0)\) rotates to \((0,1)\).
\end{examplebox}

\section{3D transformation matrices}

\subsection{\code{shift(double dx, double dy, double dz)}}
This returns
\[
\begin{bmatrix}
1 & 0 & 0 & dx \\
0 & 1 & 0 & dy \\
0 & 0 & 1 & dz \\
0 & 0 & 0 & 1
\end{bmatrix}.
\]
It translates points in \(\R^3\) by the vector \((dx,dy,dz)\).

\subsection{\code{scale(double kx, double ky, double kz)}}
This returns the diagonal homogeneous scaling matrix
\[
\begin{bmatrix}
k_x & 0 & 0 & 0 \\
0 & k_y & 0 & 0 \\
0 & 0 & k_z & 0 \\
0 & 0 & 0 & 1
\end{bmatrix}.
\]

\subsection{\code{rot_x(double angle)}}
Rotation about the \(x\)-axis is
\[
R_x(\theta) =
\begin{bmatrix}
1 & 0 & 0 & 0 \\
0 & \cos\theta & -\sin\theta & 0 \\
0 & \sin\theta & \cos\theta & 0 \\
0 & 0 & 0 & 1
\end{bmatrix}.
\]

\subsection{\code{rot_y(double angle)}}
Rotation about the \(y\)-axis is
\[
R_y(\theta) =
\begin{bmatrix}
\cos\theta & 0 & -\sin\theta & 0 \\
0 & 1 & 0 & 0 \\
\sin\theta & 0 & \cos\theta & 0 \\
0 & 0 & 0 & 1
\end{bmatrix}.
\]

\subsection{\code{rot_z(double angle)}}
Rotation about the \(z\)-axis is
\[
R_z(\theta) =
\begin{bmatrix}
\cos\theta & -\sin\theta & 0 & 0 \\
\sin\theta & \cos\theta & 0 & 0 \\
0 & 0 & 1 & 0 \\
0 & 0 & 0 & 1
\end{bmatrix}.
\]

\begin{examplebox}[Composite transform example]
Suppose a point \(p=(1,2,3)\) is represented as \((1,2,3,1)^\top\). If you first rotate around the \(z\)-axis and then translate by \((5,0,-1)\), the combined transform is
\[
T(5,0,-1)R_z(\theta).
\]
Matrix multiplication composes transformations from right to left when column vectors are used.
\end{examplebox}

\chapter{Library-Specific Utilities and Numerical Notes}

\section{\code{check_similar(matrix mat1, matrix mat2)}}
This helper is specific to ELDA. It removes all-zero rows from each matrix and then compares the remaining row sequences exactly.

\begin{techbox}[Not the standard notion of matrix similarity]
In mathematics, two square matrices \(A\) and \(B\) are called similar if there exists an invertible matrix \(P\) such that
\[
B = P^{-1}AP.
\]
ELDA's \code{check_similar()} does \emph{not} test that condition.
\end{techbox}

\begin{examplebox}[Example]
If
\[
A =
\begin{bmatrix}
1 & 2 \\
0 & 0 \\
3 & 4
\end{bmatrix},
\qquad
B =
\begin{bmatrix}
1 & 2 \\
3 & 4 \\
0 & 0
\end{bmatrix},
\]
then both reduce to the row list
\[
\{[1,2], [3,4]\},
\]
so the helper returns \code{true}.
\end{examplebox}

\section{\code{neg_zero(matrix\& m)}}
This is a display cleanup helper that replaces \(-0\) with \(0\). The mathematics is trivial because \(-0 = 0\), but floating-point output sometimes preserves the signed zero bit.

\section{Constants}
The header defines two public constants:
\begin{itemize}
    \item \code{PI = 3.141593}, used naturally in trigonometric transform examples,
    \item \code{EPS = 1e-9}, intended as a floating-point tolerance constant.
\end{itemize}
At present, \code{EPS} is more of a reserved utility constant than a value consistently used across all comparisons.

\section{Numerical caveats}
The current library uses exact comparisons against zero in several places. That is simple and readable, but in floating-point numerical linear algebra it can be fragile.

\begin{mathbox}[Practical warning]
If a pivot is extremely small rather than exactly zero, elimination without tolerance checks can amplify roundoff error. The header defines \code{EPS = 1e-9}, but the current implementation does not consistently use it for pivot tests or equality checks.
\end{mathbox}

\section{Shape prerequisites summary}
\begin{tabularx}{\linewidth}{@{}lX@{}}
\toprule
Function family & Required shape or condition \\
\midrule
\code{operator+}, \code{operator-}, \code{operator==}, \code{shape_comp} & Same dimensions \\
\code{operator*(matrix)} & Left columns = right rows \\
\code{det}, \code{minor}, \code{cofactor}, \code{adjoint}, \code{inverse} & Square matrix \\
\code{solve} & Augmented matrix of shape \(n \times (n+1)\) \\
\code{free_variable} & Current implementation expects a square coefficient matrix \\
\code{char_poly} & Mathematically intended for square matrices \\
\bottomrule
\end{tabularx}

\chapter{Engineering Structure and Maintenance Notes}

\section{Header organization}
\begin{itemize}
    \item \code{matrix.hpp} centralizes both the class interface and several free functions.
    \item \code{transforms.hpp} keeps geometry-related factories separate from algebraic routines.
    \item \code{linalg.hpp} provides one include point for consumers who want the entire API.
\end{itemize}

\section{Implementation organization}
The implementation is split by responsibility:
\begin{itemize}
    \item \code{src/matrix.cpp} handles the heavy mathematical work,
    \item \code{src/transforms.cpp} stays focused on transformation matrix construction,
    \item \code{src/vector_utils.cpp} contains the one auxiliary comparison helper.
\end{itemize}

\section{Mutability policy}
One of the main design ideas in ELDA is the distinction between:
\begin{itemize}
    \item operations that conceptually create a new matrix, such as \(A+B\), \(A^\top\), or \(A^{-1}\),
    \item operations that conceptually reduce a matrix in place, such as Gaussian elimination.
\end{itemize}
This policy keeps the user-facing code readable, but maintainers should be explicit in documentation about which functions mutate state.

\section{Starter checklist for contributors}
\begin{enumerate}[label=\arabic*.]
    \item Build the project with CMake and run the demo executable.
    \item Read \code{include/elda/matrix.hpp} and map each declaration to its implementation in \code{src/matrix.cpp}.
    \item Test shape preconditions before extending any algorithm.
    \item For numerically sensitive features, prefer tolerance-based comparisons over exact floating-point equality.
    \item Keep mathematics and naming aligned: if a helper is library-specific rather than standard, document that fact clearly.
\end{enumerate}

\appendix

\chapter{Quick Reference}

\section{Public API map}
\begin{longtable}{@{}p{0.28\textwidth}p{0.28\textwidth}p{0.34\textwidth}@{}}
\toprule
Function or helper & Main mathematical idea & Short usage note \\
\midrule
\endhead
\code{matrix()}, \code{matrix(r,c)} & Zero matrix construction & Allocates a matrix of the requested shape. \\
\code{get_element}, \code{ref_element} & Coordinate selection & Direct access to one entry. \\
\code{operator+}, \code{operator-} & Entrywise vector space operations & Requires equal shapes. \\
\code{operator*(double)} & Scalar multiplication & Multiplies every entry by the scalar. \\
\code{operator*(matrix)} & Composition of linear maps & Requires matching inner dimension. \\
\code{row_op}, \code{row_swap}, \code{row_multi} & Elementary row operations & Basis of elimination routines. \\
\code{col_op}, \code{col_swap}, \code{col_multi} & Elementary column operations & Useful for canonical forms. \\
\code{echelon}, \code{gaussian}, \code{gauss_jordan} & Elimination & Mutates the matrix in place. \\
\code{canonical} & Mixed row/column reduction & Produces a more diagonalized pattern. \\
\code{rank} & Dimension of row space & Counts nonzero rows after elimination. \\
\code{free_variable} & Rank-nullity on square systems & Returns \(n-\rankop(A)\). \\
\code{trace} & Sum of diagonal & Spectral invariant for square matrices. \\
\code{diag_product} & Product of diagonal & Equals determinant for triangular matrices. \\
\code{det} & Signed volume scaling & Computed through elimination. \\
\code{minor}, \code{cofactor}, \code{adjoint} & Cofactor expansion machinery & Defined for square matrices. \\
\code{transpose} & Reflection across the main diagonal & Swaps rows and columns. \\
\code{inverse} & Undoing a linear map & Uses Gauss-Jordan elimination. \\
\code{solve} & Solve \(Ax=b\) & Input must be an augmented matrix. \\
\code{norm}, \code{inner_product}, \code{angle} & Frobenius geometry & Treats matrices as vectors in entry space. \\
\code{matpow} & Fast repeated multiplication & Uses exponentiation by squaring. \\
\code{char_poly} & Characteristic polynomial coefficients & Uses Faddeev-LeVerrier recurrence. \\
\code{check_ortho} & Orthogonality test & Checks \(A^{-1}=A^\top\). \\
\code{check_unitary} & Legacy project helper & Not a standard unitary predicate. \\
\code{shift}, \code{scale}, \code{rotate}, \code{rot_x}, \code{rot_y}, \code{rot_z} & Homogeneous transforms & Angles are in radians. \\
\code{check_similar} & Zero-row filtered comparison & Library-specific behavior. \\
\bottomrule
\end{longtable}

\chapter{Build This Booklet}

\section{Command line}
From the repository root, run:

\begin{lstlisting}[style=eldaCpp]
make -C docs pdf
\end{lstlisting}

The generated PDF is placed in \code{docs/build/elda_booklet.pdf}.

\section{Why this format works well}
LaTeX is a strong fit for mathematical library documentation because it combines precise equations, technical prose, code listings, page structure, and stable PDF output in one source format. That makes it appropriate for a library like \ELDA{}, where understanding the theory behind each function is part of understanding the software.

\backmatter

\chapter*{Closing Note}
\addcontentsline{toc}{chapter}{Closing Note}
The strongest documentation for a numerical library is documentation that respects both mathematics and code. ELDA already has a coherent educational core: matrices, elimination, determinants, inverses, and transforms. This booklet formalizes that core into a reference that can support further work, testing, refactoring, and extension.

\end{document}
